% GENERATED FILE. Edit canonical lessons, then run npm run generate:books.
% canonical-source-hash: fnv1a64:6c5cc376d8145c6c
% canonical-lessons: JA-C170-saraishuu, JA-C170-sarainen, JA-C170-shougatsu, JA-C170-iwau, JA-C170-sugosu

\chapter{The New Year and the Weeks After Next}
\label{ch:ja-the-new-year-and-the-weeks-after-next}
\hlchaptermodality{170}

\begin{chapteropening}
\textbf{By the end of this chapter:} \emph{I can talk about holidays and plans further ahead.}

Complete the last lesson of chapter 170: I can talk about holidays and plans further ahead.
\end{chapteropening}

\section[\texorpdfstring{\ja{さらいしゅう} (saraishū)}{さらいしゅう (saraishū)}]{\ja{さらいしゅう} (saraishū) — the week after next}

\label{lesson:JA-C170-saraishuu}



\begin{quote}
Before the new one: say the Japanese for every month, then the Japanese for to repeat.
\end{quote}

\subsection*{You'll want to know: \ja{さらいしゅう}}
\textbf{\ja{さらいしゅう}} — \emph{saraishū} — \textquotedblleft{}the week after next\textquotedblright{}.

\begin{quote}
\textbf{\ja{さらいしゅう} \ja{ひっこします。}} — \emph{saraishū hikkoshimasu} — \textquotedblleft{}I will move house the week after next.\textquotedblright{}
\end{quote}

\begin{cousinweb}
\textbf{\ja{さらいしゅう}} (\emph{saraishū}) is Sino-Japanese (kango): three Chinese characters read with their on'yomi, readings borrowed from Chinese. They are \emph{sa}, \textquotedblleft{}again\textquotedblright{} (Middle Chinese \emph{tsojH}), \emph{rai}, \textquotedblleft{}coming\textquotedblright{} (\emph{loj}) and \emph{shū}, \textquotedblleft{}week\textquotedblright{}, whose Middle Chinese reading the sources here do not give (KANJIDIC2; Middle Chinese in Baxter's transcription).
\end{cousinweb}

\begin{grammarlens}[title={What you've built}]
One more word for everyday life at A2.
\end{grammarlens}

\subsection*{Your turn}
\begin{itemize}
\raggedright
  \item \emph{Say it:} \emph{saraishū}
  \item \emph{Say it:} \emph{saraishū}, once more
  \item \emph{Say it:} \emph{saraishū hikkoshimasu} — I will move house the week after next.
  \item \emph{Recall:} say the Japanese for every month, then the Japanese for to repeat, then say \emph{saraishū} again
\end{itemize}

\begin{tcolorbox}[breakable,colback=teal!4,colframe=teal!35!black,title={Before you move on}]
What does \ja{さらいしゅう} mean? (\textquotedblleft{}The week after next\textquotedblright{}.) Say it once more.
\end{tcolorbox}
\section[\texorpdfstring{\ja{さらいねん} (sarainen)}{さらいねん (sarainen)}]{\ja{さらいねん} (sarainen) — the year after next}

\label{lesson:JA-C170-sarainen}



\begin{quote}
Before the new one: say the Japanese for to repeat, then the Japanese for the week after next.
\end{quote}

\subsection*{You'll want to know: \ja{さらいねん}}
\textbf{\ja{さらいねん}} — \emph{sarainen} — \textquotedblleft{}the year after next\textquotedblright{}.

\begin{quote}
\textbf{\ja{さらいねん} \ja{いなかに} \ja{すむ} \ja{つもりです。}} — \emph{sarainen inaka ni sumu tsumori desu} — \textquotedblleft{}I plan to live in the countryside the year after next.\textquotedblright{}
\end{quote}

\begin{cousinweb}
\textbf{\ja{さらいねん}} (\emph{sarainen}) is Sino-Japanese (kango): three Chinese characters read with their on'yomi, readings borrowed from Chinese. They are \emph{sa}, \textquotedblleft{}again\textquotedblright{} (Middle Chinese \emph{tsojH}), \emph{rai}, \textquotedblleft{}coming\textquotedblright{} (\emph{loj}) and \emph{nen}, \textquotedblleft{}year\textquotedblright{} (\emph{nen}) (KANJIDIC2; Middle Chinese in Baxter's transcription).
\end{cousinweb}

\begin{grammarlens}[title={What you've built}]
One more word for everyday life at A2.
\end{grammarlens}

\subsection*{Your turn}
\begin{itemize}
\raggedright
  \item \emph{Say it:} \emph{sarainen}
  \item \emph{Say it:} \emph{sarainen}, once more
  \item \emph{Say it:} \emph{sarainen inaka ni sumu tsumori desu} — I plan to live in the countryside the year after next.
  \item \emph{Recall:} say the Japanese for to repeat, then the Japanese for the week after next, then say \emph{sarainen} again
\end{itemize}

\begin{tcolorbox}[breakable,colback=teal!4,colframe=teal!35!black,title={Before you move on}]
What does \ja{さらいねん} mean? (\textquotedblleft{}The year after next\textquotedblright{}.) Say it once more.
\end{tcolorbox}
\section[\texorpdfstring{\ja{しょうがつ} (shōgatsu)}{しょうがつ (shōgatsu)}]{\ja{しょうがつ} (shōgatsu) — the New Year; January}

\label{lesson:JA-C170-shougatsu}



\begin{quote}
Before the new one: say the Japanese for the week after next, then the Japanese for the year after next.
\end{quote}

\subsection*{You'll want to know: \ja{しょうがつ}}
\textbf{\ja{しょうがつ}} — \emph{shōgatsu} — \textquotedblleft{}the New Year; January\textquotedblright{}.

\begin{quote}
\textbf{\ja{しょうがつに} \ja{いなかに} \ja{かえります。}} — \emph{shōgatsu ni inaka ni kaerimasu} — \textquotedblleft{}I will go back to the countryside at New Year.\textquotedblright{}
\end{quote}

\begin{cousinweb}
\textbf{\ja{しょうがつ}} (\emph{shōgatsu}) is Sino-Japanese (kango): two Chinese characters read with their on'yomi, readings borrowed from Chinese. They are \emph{shō}, \textquotedblleft{}correct, righteous\textquotedblright{} (Middle Chinese \emph{tsyeng}) and \emph{gatsu}, \textquotedblleft{}moon, month\textquotedblright{} (\emph{ngjwot}) (KANJIDIC2; Middle Chinese in Baxter's transcription).
\end{cousinweb}

\begin{grammarlens}[title={What you've built}]
One more word for everyday life at A2.
\end{grammarlens}

\subsection*{Your turn}
\begin{itemize}
\raggedright
  \item \emph{Say it:} \emph{shōgatsu}
  \item \emph{Say it:} \emph{shōgatsu}, once more
  \item \emph{Say it:} \emph{shōgatsu ni inaka ni kaerimasu} — I will go back to the countryside at New Year.
  \item \emph{Recall:} say the Japanese for the week after next, then the Japanese for the year after next, then say \emph{shōgatsu} again
\end{itemize}

\begin{tcolorbox}[breakable,colback=teal!4,colframe=teal!35!black,title={Before you move on}]
What does \ja{しょうがつ} mean? (\textquotedblleft{}The New Year; January\textquotedblright{}.) Say it once more.
\end{tcolorbox}
\section[\texorpdfstring{\ja{いわう} (iwau)}{いわう (iwau)}]{\ja{いわう} (iwau) — to celebrate}

\label{lesson:JA-C170-iwau}



\begin{quote}
Before the new one: say the Japanese for the year after next, then the Japanese for the New Year.
\end{quote}

\subsection*{You'll want to know: \ja{いわう}}
\textbf{\ja{いわう}} — \emph{iwau} — \textquotedblleft{}to celebrate\textquotedblright{}.

\begin{quote}
\textbf{\ja{らいしゅう} \ja{ちちの} \ja{たんじょうびを} \ja{いわいます。}} — \emph{raishū chichi no tanjōbi o iwaimasu} — \textquotedblleft{}Next week we will celebrate my father's birthday.\textquotedblright{}
\end{quote}

Polite \textbf{\ja{いわいます}}, \emph{iwaimasu}; a godan verb.

\begin{cousinweb}
\textbf{\ja{いわう}} (\emph{iwau}) is a native Japanese word (wago), inherited rather than borrowed from Chinese; its kanji was attached later as a kun'yomi, and its older origin is not settled (JMdict, KANJIDIC2).
\end{cousinweb}

\begin{grammarlens}[title={What you've built}]
One more word for everyday life at A2.
\end{grammarlens}

\subsection*{Your turn}
\begin{itemize}
\raggedright
  \item \emph{Say it:} \emph{iwau}
  \item \emph{Say it:} \emph{iwau}, once more
  \item \emph{Say it:} \emph{raishū chichi no tanjōbi o iwaimasu} — Next week we will celebrate my father's birthday.
  \item \emph{Recall:} say the Japanese for the year after next, then the Japanese for the New Year, then say \emph{iwau} again
\end{itemize}

\begin{tcolorbox}[breakable,colback=teal!4,colframe=teal!35!black,title={Before you move on}]
What does \ja{いわう} mean? (\textquotedblleft{}To celebrate\textquotedblright{}.) Say it once more.
\end{tcolorbox}
\section[\texorpdfstring{\ja{すごす} (sugosu)}{すごす (sugosu)}]{\ja{すごす} (sugosu) — to spend time; to pass (the days)}

\label{lesson:JA-C170-sugosu}



\begin{quote}
Before the new one: say the Japanese for the New Year, then the Japanese for to celebrate.
\end{quote}

\subsection*{You'll want to know: \ja{すごす}}
\textbf{\ja{すごす}} — \emph{sugosu} — \textquotedblleft{}to spend time; to pass (the days)\textquotedblright{}.

\begin{quote}
\textbf{\ja{しょうがつは} \ja{かぞくと} \ja{すごします。}} — \emph{shōgatsu wa kazoku to sugoshimasu} — \textquotedblleft{}I will spend the New Year with my family.\textquotedblright{}
\end{quote}

Polite \textbf{\ja{すごします}}, \emph{sugoshimasu}; a godan verb.

\begin{cousinweb}
\textbf{\ja{すごす}} (\emph{sugosu}) is a native Japanese word (wago), inherited rather than borrowed from Chinese; its kanji was attached later as a kun'yomi, and its older origin is not settled (JMdict, KANJIDIC2).
\end{cousinweb}

\begin{grammarlens}[title={What you've built}]
One more word for everyday life at A2.
\end{grammarlens}

\subsection*{Your turn}
\begin{itemize}
\raggedright
  \item \emph{Say it:} \emph{sugosu}
  \item \emph{Say it:} \emph{sugosu}, once more
  \item \emph{Say it:} \emph{shōgatsu wa kazoku to sugoshimasu} — I will spend the New Year with my family.
  \item \emph{Recall:} say the Japanese for the New Year, then the Japanese for to celebrate, then say \emph{sugosu} again
\end{itemize}

\begin{tcolorbox}[breakable,colback=teal!4,colframe=teal!35!black,title={Before you move on}]
What does \ja{すごす} mean? (\textquotedblleft{}To spend time; to pass (the days)\textquotedblright{}.) Say it once more.
\end{tcolorbox}
